\documentclass[10pt]{article}
\usepackage[a4paper,margin=2.2cm]{geometry}
\usepackage[T1]{fontenc}
\usepackage{lmodern}
\usepackage{amsmath,amssymb}
\usepackage{booktabs}
\usepackage{graphicx}
\usepackage{caption}
\usepackage{float}
\usepackage{xcolor}
\usepackage{enumitem}
\usepackage[hidelinks]{hyperref}
\captionsetup{font=small,labelfont=bf}
\setlist{nosep,leftmargin=*}

\newcommand{\LiveRuns}{15}
\newcommand{\LiveBlocksObserved}{183}
\newcommand{\LiveReorgBlocks}{43}
\newcommand{\LiveReorgPct}{23.5}
\newcommand{\LiveMaxReorg}{5}
\newcommand{\FinalityDepth}{6}
\newcommand{\TamperCases}{143,000}
\newcommand{\TDR}{100.00}
\newcommand{\TDRlo}{99.997}
\newcommand{\FPR}{0.00}
\newcommand{\FPRhi}{0.014}
\newcommand{\BenignCases}{26,000}
\newcommand{\InsSOneEmu}{100.0}
\newcommand{\InsSOneBesu}{100.0}
\newcommand{\InsSOneDB}{100.0}
\newcommand{\InsSOnePos}{100.0}
\newcommand{\InsSTwoEmu}{100.0}
\newcommand{\InsSTwoBesu}{100.0}
\newcommand{\InsSTwoDB}{0.0}
\newcommand{\InsSTwoPos}{100.0}
\newcommand{\InsSThreeEmu}{100.0}
\newcommand{\InsSThreeBesu}{100.0}
\newcommand{\InsSThreeDB}{0.0}
\newcommand{\InsSThreePos}{100.0}
\newcommand{\MaxTpsEmu}{187}
\newcommand{\TpsTwoHEmu}{184}
\newcommand{\LatTwoHEmu}{1.99}
\newcommand{\MaxTpsBesu}{57}
\newcommand{\TpsTwoHBesu}{58}
\newcommand{\LatTwoHBesu}{46.21}
\newcommand{\MaxTpsDB}{1,378}
\newcommand{\TpsTwoHDB}{200}
\newcommand{\LatTwoHDB}{0.00}
\newcommand{\MaxTpsRatio}{3.3}
\newcommand{\InclWtwoEmu}{1.56}
\newcommand{\FinalWtwoEmu}{7.55}
\newcommand{\TpsWtwoEmu}{49.6}
\newcommand{\InclWtwoBesu}{1.28}
\newcommand{\FinalWtwoBesu}{1.28}
\newcommand{\TpsWtwoBesu}{49.7}
\newcommand{\InclWtwoDB}{0.00}
\newcommand{\FinalWtwoDB}{0.00}
\newcommand{\TpsWtwoDB}{50.0}
\newcommand{\InclWtwoPos}{213.21}
\newcommand{\FinalWtwoPos}{1,103}
\newcommand{\TpsWtwoPos}{16.1}
\newcommand{\FinalRatioPos}{146}
\newcommand{\InclRatioPos}{137}
\newcommand{\FeePos}{2.30}
\newcommand{\CostEmu}{35.9}
\newcommand{\EdynEmu}{0.050}
\newcommand{\EattEmu}{2.03}
\newcommand{\CostBesu}{35.9}
\newcommand{\EdynBesu}{0.302}
\newcommand{\EattBesu}{2.03}
\newcommand{\CostDB}{25.7}
\newcommand{\EdynDB}{0.008}
\newcommand{\EattDB}{0.51}
\newcommand{\CostPos}{2,300,688}
\newcommand{\EdynPos}{0.171}
\newcommand{\EattPos}{6.25}
\newcommand{\EdynRatioBesu}{6.0}
\newcommand{\EattRatioPos}{3.1}
\newcommand{\VerifyEmu}{14.2}
\newcommand{\VerifyBesu}{33.0}
\newcommand{\VerifyDB}{5.1}
\newcommand{\VTpsEmuFour}{177}
\newcommand{\VTpsBesuFour}{60}
\newcommand{\VTpsEmuSeven}{175}
\newcommand{\VTpsBesuSeven}{54}
\newcommand{\VTpsEmuTen}{175}
\newcommand{\VTpsBesuTen}{--}
\newcommand{\BesuTenSuccess}{4}
\newcommand{\GasKone}{153,376}
\newcommand{\GasKhundred}{2,390}
\newcommand{\PosTpsKone}{15.4}
\newcommand{\PosTpsKfifty}{44.9}
\newcommand{\PosFeeKfifty}{0.051}
\newcommand{\KWminEmu}{0.15}
\newcommand{\Trials}{777}
\newcommand{\LiveValidators}{4}
\newcommand{\LiveRTT}{15.7}
\newcommand{\LiveGasLimit}{30,000,000}
\newcommand{\LiveInTurn}{71}

\newcommand{\poa}{PoA\textsuperscript{2}}
\newcommand{\figif}[3]{\IfFileExists{figures/#1.pdf}{%
  \begin{figure}[t]\centering\includegraphics[width=#2]{figures/#1.pdf}\caption{#3}\label{fig:#1}\end{figure}}{}}

\title{FinVerity on PureChain: Experimental Setup and Results\\[2pt]
\large Blockchain-anchored, continuously verifiable financial audit evidence}
\author{}
\date{Generated from the complete measurement campaign (\Trials{} trials)}

\begin{document}
\maketitle

\begin{abstract}
\noindent This report documents the experimental methodology and the measured results for FinVerity, an audit-evidence
architecture that commits SHA-256 digests of canonicalised financial evidence packages to the PureChain
\poa{} blockchain while the evidence itself remains off-chain. Four transaction-recording infrastructures are
compared under identical workloads derived from public synthetic financial transaction datasets: a centralised
database (PostgreSQL), a conventional permissioned blockchain (Hyperledger Besu QBFT), a conventional public-chain
configuration (Ethereum proof-of-stake with mainnet timing) and PureChain, measured both on the live network and on a
\poa{} emulation calibrated to it. The evaluation covers performance (latency, finality, throughput), efficiency
(energy, cost, resources) and auditability (verification, tamper detection, insider resistance), and every core
experiment is replicated on three independent transaction datasets (MoMTSim, SAML-D, IBM AML). All tables and
figures are generated directly from the raw measurements; Table~\ref{tab:status} states which experiments are
complete.
\end{abstract}

% =====================================================================================
\section{Experimental Setup}\label{sec:setup}

\subsection{Research design}
The experiments address the question whether a lightweight, zero-gas, \poa-based permissioned blockchain can provide
cryptographically verifiable financial audit evidence while reducing transaction latency, energy consumption and
operational cost. Equivalent workloads are executed on every infrastructure and three groups of metrics are
collected: \emph{performance} (latency, finality, throughput), \emph{efficiency} (energy, cost, CPU, memory, network,
storage) and \emph{auditability} (verification success, evidence--transaction binding, tamper detection, insider
resistance).

\begin{table}[H]
\centering\scriptsize
\caption{Evaluation framework: metric groups and their operationalisation.}\label{tab:metrics}
\begin{tabular}{p{2.1cm}p{4.3cm}p{8.6cm}}
\toprule
Group & Metric & Operationalisation \\
\midrule
Performance & Throughput (TPS) & committed records / (last inclusion $-$ first arrival) \\
 & Transaction latency & arrival $\rightarrow$ first observation of the canonical block containing the record \\
 & Finality latency & arrival $\rightarrow$ finality under the rule of Table~\ref{tab:infra} \\
 & Confirmation time & arrival $\rightarrow$ observation of block $b+k$, $k\in\{1,2,3,6\}$ \\
\midrule
Efficiency & Energy per transaction (J/tx) & attributed node energy and measured dynamic CPU energy per record \\
 & CPU utilisation & CPU-seconds of client and infrastructure containers per second of processing \\
 & Memory utilisation & peak resident memory of infrastructure containers \\
 & Network I/O & request and response bytes per record \\
 & Storage overhead & on-chain bytes per record; off-chain evidence bytes per record \\
 & Operational cost per transaction & $C_{\mathit{compute}}+C_{\mathit{storage}}+C_{\mathit{network}}+C_{\mathit{fee}}$ \\
\midrule
Auditability & Evidence verification time & record lookup + evidence retrieval + re-hash, per record \\
 & Verification success rate & share of committed records with $H_i'=H_i$ \\
 & Tamper detection rate & TDR over eleven modification classes; FPR over benign transformations \\
 & Evidence--transaction binding & identifier, commitment and approver all match \\
 & Insider resistance & attacks by a privileged operator prevented or detected \\
\bottomrule
\end{tabular}
\end{table}

\subsection{Infrastructures under test}\label{sec:infra}
\begin{table}[H]
\centering\scriptsize
\caption{Infrastructures compared. All blockchain configurations run the same Solidity contract
(\texttt{FinVerityRegistry}, solc 0.8.19, EVM target \emph{paris}).}\label{tab:infra}
\begin{tabular}{p{2.6cm}p{3.5cm}p{3.1cm}p{2.4cm}p{3.3cm}}
\toprule
Configuration & Consensus / engine & Nodes & Finality rule & Fees \\
\midrule
PureChain (live) & \poa{} (Geth 1.13.15 fork, Clique-style seals) & \LiveValidators{} validators (remote, public RPC) & \FinalityDepth{} confirmations$^{a}$ & gas price 0 \\
\poa{} emulation & Clique PoA, Geth 1.13.15, 1\,s period & 4 signers (Docker, 1 CPU each) & \FinalityDepth{} confirmations$^{a}$ & fee-free (1\,wei nominal)$^{b}$ \\
Besu QBFT & QBFT BFT, Besu 24.12.2, 1\,s period & 4 validators (Docker, 1 CPU each) & immediate (inclusion) & gas price 0 \\
Ethereum PoS & Casper FFG + LMD-GHOST; Geth 1.13.15 + Prysm 5.0.4 & 1 node, 64 validators & finalized checkpoint & EIP-1559$^{c}$ \\
PostgreSQL 16 & single operator, \texttt{fsync=on}, synchronous commit & 1 server (Docker) & durable commit & none \\
\bottomrule
\end{tabular}
\par\smallskip\footnotesize
$^{a}$Forks up to depth \LiveMaxReorg{} were observed on live PureChain (Section~\ref{sec:live});
$k=\FinalityDepth$ is the smallest depth exceeding every observed reorganisation. $^{b}$Upstream Geth enforces a 1\,wei floor for transactions relayed between
peers; the cost is economically zero ($\approx1.5\times10^{-13}$\,ETH per record) and no fee is charged in the cost
model. $^{c}$Fees are expressed at a mainnet reference price of 5\,gwei and 3{,}000\,USD/ETH.
\end{table}

Every blockchain configuration uses a 30\,M block gas limit. Besu and the \poa{} emulation give every validator the
same 1-CPU quota, so that validator-count scaling is not confounded by contention on the single test host. The
Ethereum PoS configuration uses mainnet timing (12\,s slots, 32 slots per epoch) and represents the
\emph{conventional public-chain configuration} named in the hypothesis. Only the backend under test runs during a
trial; all other local networks are stopped.

\subsection{Testbed}
Client and local infrastructure run on one workstation: Intel Core i5-12400F (6 cores, 12 threads, 65\,W base
power), 32\,GB RAM, Windows 11, Docker Desktop 29.4 (12 CPUs, 16.7\,GB assigned). Software: Python 3.13,
\texttt{web3.py} 6.20.4, \texttt{purechainlib} 2.1.7, PostgreSQL 16, Hyperledger Besu 24.12.2, Geth 1.13.15,
Prysm 5.0.4. Live PureChain is reached over the public RPC endpoint (chain ID 900520900520, median RPC round trip
\LiveRTT\,ms). The measured network parameters of PureChain (\LiveValidators{} validators sealing round-robin,
\LiveInTurn\% of blocks sealed in turn, \LiveGasLimit{} block gas limit, gas price 0) were obtained by recovering the
signer of each block from its Clique-style seal.

\subsection{Financial transaction workloads}
The datasets are used strictly as \emph{financial transaction workloads}; fraud and laundering labels are discarded.
\begin{table}[H]
\centering\scriptsize
\caption{Transaction workload datasets.}\label{tab:data}
\begin{tabular}{p{4.3cm}rllr}
\toprule
Dataset & Transactions & Period & Main types & Median amount \\
\midrule
MoMTSim (primary; Azamuke et al., 2024) & 4{,}225{,}958 & 8 simulated days & transfer, payment & 596.69 \\
SAML-D (Oztas et al., 2023) & 9{,}504{,}852 & Oct 2022--Aug 2023 & card, cheque, ACH & 6{,}113.72 \\
IBM AML HI-Small (Altman et al., 2023) & 5{,}078{,}345 & Sep 2022 & cheque, card, ACH & 1{,}414.54 \\
\bottomrule
\end{tabular}
\end{table}
\noindent Using three datasets with different transaction mixes, currencies, amount distributions and account
structures tests whether the results depend on the workload source. MoMTSim drives all experiments; the core
workloads W1, W2 and W5, the tamper-detection experiment and the insider experiment are replicated on SAML-D and IBM
AML (E6--E8). All datasets are mapped to a common schema (identifier, timestamp, type, amount as a decimal string,
currency, originator, beneficiary, balances, channel). MoMTSim provides neither identifiers nor wall-clock times; both
are derived deterministically from the row index and the simulated hour. Each trial draws a contiguous chronological
window at a seeded random offset, which preserves temporal ordering and account interactions.

\subsection{Evidence generation and cryptographic commitment}
For every transaction $T_i$ a synthetic audit-evidence package $E_i=\{M_i,D_i,A_i,L_i,C_i\}$ is generated
deterministically: metadata $M_i$; invoice and purchase-order record $D_i$ with supporting documents; authorisation
record $A_i$ (approval tier by amount, approver, segregation of duties); balanced double-entry journal $L_i$; and
context $C_i$. Complexity levels are L1 (metadata), L2 (+invoice), L3 (+authorisation) and L4 (complete).
Supporting documents are incompressible byte streams of controlled size (0, 10\,KB, 100\,KB, 1\,MB, 5\,MB), each
represented in the package by its SHA-256 digest. The commitment is
\begin{equation}
E_i^{*}=\operatorname{Canon}(E_i),\qquad H_i=\mathrm{SHA256}(E_i^{*}),
\end{equation}
where $\operatorname{Canon}$ is an RFC~8785 profile (sorted keys, no insignificant whitespace, UTF-8, Unicode NFC,
monetary values as decimal strings, floats rejected). An approver signs (EIP-191) the domain-bound digest
\begin{equation}
S_i=\operatorname{Sign}_{\mathit{approver}}\bigl(\mathrm{keccak}(\mathit{chainId}\,\|\,\mathit{registry}\,\|\,ID_i\,\|\,H_i\,\|\,\mathit{metaHash}\,\|\,\mathit{level})\bigr).
\end{equation}
The contract accepts \texttt{recordEvidence} only from registered submitters with a valid approver signature, rejects
duplicate identifiers and stores $\{H_i,\mathit{metaHash},\mathit{submitter},\mathit{approver},t_i,\mathit{block}\}$.
In batched mode (\texttt{recordBatch}) $k$ records are committed as one Merkle root over leaves
$\mathrm{keccak}(ID_i\,\|\,H_i)$; each record keeps an inclusion proof, verifiable off-chain or with the
contract's \texttt{verifyInclusion}. The evidence is kept off-chain in a durable store (SQLite, WAL,
\texttt{synchronous=FULL}).

\subsection{Submission and measurement}
Records arrive as an open-loop Poisson process at the offered rate. Transactions are signed in advance and submitted
by 16 submitter accounts in parallel over persistent keep-alive connections. A batch closes when it holds $k$ records
or its first record has waited 1\,s. A fork-aware block watcher polls the chain head every 50\,ms (250\,ms for PoS),
records every block version with its first-observation time and walks back along parent hashes, so replaced blocks
are captured. Final transaction status is taken from receipts re-read against the canonical chain. Latencies are
measured per record from its arrival:
\begin{itemize}
\item \emph{inclusion latency}: arrival to first observation of the canonical block containing the record;
\item \emph{$k$-confirmation latency}: arrival to observation of block $b+k$;
\item \emph{finality latency}: arrival to finality under the rule of Table~\ref{tab:infra};
\item \emph{throughput}: committed records divided by the time from the first arrival to the last inclusion.
\end{itemize}

\subsection{Efficiency metrics}
Energy is reported in two ways for every backend. The \emph{attributed} energy
$E_{\mathit{tx}}=\bigl(n_{\mathit{nodes}}\,P_{\mathit{node}}\,T_{\mathit{proc}}+E_{\mathit{client}}\bigr)/N$
charges every infrastructure node ($P_{\mathit{node}}=25$\,W, $n_{\mathit{nodes}}=4$ for blockchains, 1 for the
database) for the processing window $T_{\mathit{proc}}$ from the first arrival to the last inclusion, so waiting for
finality is not charged. The \emph{measured dynamic} energy converts CPU time of the client and of the local containers
into energy with a linear power model $E=P_{\mathit{dyn}}\cdot t_{\mathit{CPU}}/n_{\mathit{cores}}$ (Fan et al., 2007;
$P_{\mathit{dyn}}=65$\,W, 12 threads). Cost is
$C_{\mathit{tx}}=C_{\mathit{compute}}+C_{\mathit{storage}}+C_{\mathit{network}}+C_{\mathit{fee}}$, with 0.12\,USD/kWh,
0.023\,USD/GB-month off-chain storage, 0.08\,USD/GB-month on-chain storage per replica, 0.09\,USD/GB network traffic
and an 84-month retention period. $C_{\mathit{fee}}$ is non-zero only for the fee-bearing public-chain configuration.

\subsection{Auditability metrics}
The auditor retrieves $H_i$ (or the batch root and proof), recomputes $H_i'=\mathrm{SHA256}(\operatorname{Canon}(E_i'))$
after re-hashing every stored document, and sets $\operatorname{Verify}(E_i')=1$ iff $H_i'=H_i$. Binding success also
requires the identifier and the recorded approver to match the package. The tamper detection rate is
\begin{equation}
\mathrm{TDR}=\frac{\text{detected modifications}}{\text{applicable modifications}}\times100,
\end{equation}
over eleven malicious modification classes (amount, date, invoice number, vendor, authorisation, supporting document,
deleted document, ledger entry, beneficiary, single-bit change, record substitution). Two benign transformations
(re-serialisation with reordered keys, Unicode NFD) give the false-positive rate. The insider experiment assumes an
operator holding submitter and approver credentials who (S1) alters evidence only, (S2) alters evidence and overwrites
the stored commitment, or (S3) fabricates and backdates a replacement record.

\subsection{Experimental design}
\begin{table}[H]
\centering\scriptsize
\caption{Experiments (all use evidence level L4 unless varied).}\label{tab:design}
\begin{tabular}{llp{9.3cm}}
\toprule
ID & Factor & Configuration \\
\midrule
E1 & workload & W1 1k/10\,KB ($n$=30), W2 10k/10\,KB (10), W3 50k/100\,KB (5), W4 100k/100\,KB (3), W5 10k/1\,MB (10); 50 rec/s \\
E2 & offered load & 10, 25, 50, 100, 200, 400 rec/s and unthrottled; 5 trials \\
E3 & evidence size & none, 10\,KB, 100\,KB, 1\,MB, 5\,MB; 1{,}000 records; 5 trials \\
E4 & complexity & L1--L4; 1{,}000 records, 100\,KB; 5 trials \\
E5 & validators & 4, 7, 10 (Besu and \poa{} emulation) at 50 and 200 rec/s; 5 trials \\
E6 & tamper detection & 1{,}000 records per dataset, 13 transformations per record; 3 trials per dataset \\
E7 & dataset & W1 ($n$=10), W2 (5) and W5 (3) replicated on SAML-D and IBM AML \\
E8 & insider & 300 records, 100 attacked per scenario; 3 trials (MoMTSim), 2 per secondary dataset \\
E9 & batch size & $k\in\{1,10,25,50,100\}$ at 50 and 200 rec/s; 5 trials \\
\bottomrule
\end{tabular}
\end{table}

\subsection{Statistical analysis}
Each configuration is repeated in independent trials with different seeded workload windows. Trial-level results are
reported as mean with a 95\% $t$-interval; latency distributions as median, P95 and P99. Backends are compared with
Welch's $t$-test and the Mann--Whitney $U$ test with Holm--Bonferroni correction and Cliff's $\delta$ as effect size.
Detection rates are reported with exact Clopper--Pearson intervals.

\subsection{\poa{} emulation and its calibration}\label{sec:emulation}
The public PureChain RPC endpoint stopped producing blocks during the campaign (Section~\ref{sec:validity}). To
complete the controlled experiments, a \poa{} emulation was built from the parameters measured on the live network and
validated against a live run with an identical workload (Table~\ref{tab:calibration}). Block packing and
confirmation-depth latency agree closely; inclusion is faster on the emulation because the wide-area path and the
public load balancer are absent. The emulation did not reproduce the forks seen on the live network under heavy load;
applying the same \FinalityDepth-confirmation rule to both keeps the emulated finality latency conservative. Results
from the emulation are labelled as such throughout.
\begin{table}[t]
\centering
\small
\caption{Calibration of the PoA\textsuperscript{2} emulation against live PureChain (identical workload: 300 MoMTSim records, 10\,KB evidence, 50 records/s Poisson arrivals).}
\label{tab:calibration}
\begin{tabular}{lrrr}
\toprule
Metric & Live PureChain & PoA\textsuperscript{2} emulation & Deviation \\
\midrule
Records committed & 300 & 300 & +0.0\% \\
Throughput (records/s) & 36.7 & 40.0 & +8.9\% \\
Inclusion latency, median (s) & 2.24 & 1.49 & -33.2\% \\
Inclusion latency, P95 (s) & 3.22 & 1.96 & -39.2\% \\
1-confirmation latency, median (s) & 3.07 & 2.48 & -19.2\% \\
3-confirmation latency, median (s) & 4.29 & 4.48 & +4.4\% \\
6-confirmation latency, median (s) & 7.19 & 7.47 & +3.8\% \\
Transactions per non-empty block, mean & 42.9 & 42.9 & +0.0\% \\
Distinct block signers & 4 & 4 & +0.0\% \\
Gas per record & 153,317 & 153,434 & +0.1\% \\
Verification success (\%) & 100.0 & 100.0 &  \\
\bottomrule
\end{tabular}
\par\smallskip\footnotesize Live inclusion latency additionally contains the wide-area round trip and the public load balancer.
\end{table}


% =====================================================================================
\section{Results}\label{sec:results}
\begin{table}[t]
\centering
\scriptsize
\caption{Campaign completion status (completed / planned trials) at the time this report was generated. PureChain (live) trials are pending the recovery of the public RPC endpoint.}
\label{tab:status}
\begin{tabular}{llccccc}
\toprule
ID & Experiment & PureChain (live) & PoA\textsuperscript{2} emulation & Besu QBFT & PostgreSQL & Ethereum PoS \\
\midrule
E1 & Controlled workloads W1--W5 & 0/58 & 58/58 & 58/58 & 58/58 & 7/7 \\
E2 & Offered-load sweep & 0/35 & 35/35 & 35/35 & 35/35 & 4/4 \\
E3 & Evidence-size sweep & 0/25 & 25/25 & 25/25 & 25/25 &  \\
E4 & Evidence-complexity sweep & 0/20 & 20/20 & 20/20 & 20/20 &  \\
E5 & Validator-count sweep &  & 30/30 & 24/30 &  &  \\
E6 & Tamper detection & 4/9 & 9/9 &  &  &  \\
E7 & Cross-dataset replication & 0/36 & 36/36 & 36/36 & 36/36 &  \\
E8 & Insider tampering & 0/7 & 7/7 & 7/7 & 7/7 & 2/2 \\
E9 & Batched Merkle commitments & 0/50 & 50/50 & 50/50 & 50/50 & 4/4 \\
\bottomrule
\end{tabular}
\end{table}


\subsection{Live PureChain characterisation}\label{sec:live}
Block seals recovered from the live chain show \LiveValidators{} authorised validators sealing in round-robin, with
in-turn (difficulty 2) and out-of-turn (difficulty 1) blocks, which is Clique-style proof-of-authority behaviour. Blocks
are produced at about 1\,s intervals while transactions are pending and the chain idles otherwise. Under sustained
load, competing blocks at the same height were produced by different signers: over \LiveRuns{} live trials,
\LiveReorgBlocks{} of \LiveBlocksObserved{} observed blocks (\LiveReorgPct\%) were replaced by a competing block, with
forks up to \LiveMaxReorg{} blocks deep. Every transaction contained in a replaced block was re-included and verified.
Inclusion therefore gives only probabilistic finality, and \FinalityDepth{} confirmations are used as the finality
rule for \poa{} (live and emulated). Shortly after the heaviest load the public network stopped producing blocks
(Section~\ref{sec:validity}).
\begin{table}[t]
\centering
\scriptsize
\caption{Measurements on the live PureChain network (4 validators) before the public RPC became unavailable (E6 runs, 1{,}000 records each, 100\,KB evidence, 50 records/s offered).}
\label{tab:live}
\resizebox{\linewidth}{!}{%
\begin{tabular}{lrrrrrrrr}
\toprule
Dataset & Committed & rec/s & Incl. P50 (s) & 3-conf P50 (s) & tx/block & Reorged & Depth & Verified (\%) \\
\midrule
momtsim & 1,000 & 36.4 & 4.30 & 10.75 & 66.7 & 2 & 1 & 100.0 \\
momtsim & 1,000 & 41.2 & 5.14 & 12.64 & 90.9 & 7 & 3 & 100.0 \\
momtsim & 1,000 & 27.4 & 10.37 & 20.70 & 76.9 & 5 & 3 & 100.0 \\
saml\_d & 1,000 & 34.7 & 6.86 & 13.92 & 83.3 & 5 & 3 & 100.0 \\
\bottomrule
\end{tabular}}
\par\smallskip\footnotesize Submitted through the load-balanced endpoint; nonce splitting across two RPC nodes inflates latency (see Section~\ref{sec:validity}).
\end{table}


\subsection{Evidence integrity (E6)}
Across \TamperCases{} applicable malicious modifications, the verification detected \TDR\% (95\% CI lower bound
\TDRlo\%), while \BenignCases{} semantically neutral transformations produced a false-positive rate of \FPR\%
(upper bound \FPRhi\%). Canonicalisation therefore makes the commitment insensitive to serialisation while any change
to content, documents or authorisation invalidates it (Table~\ref{tab:tamper}).
\begin{table}[t]
\centering
\scriptsize
\caption{Tamper detection against on-chain commitments (E6). Malicious rows report the detection rate (TDR); benign rows report the false-positive rate (FPR). Exact Clopper--Pearson 95\% intervals. Pooled over PureChain (live): 4 runs; PoA\textsuperscript{2} emulation: 9 runs.}
\label{tab:tamper}
\begin{tabular}{llrrrc}
\toprule
Kind & Modification & Cases & Flagged & Rate (\%) & 95\% CI \\
\midrule
benign & reserialize\_reorder & 13,000 & 0 & 0.00 & [0.000, 0.028] \\
benign & unicode\_nfd & 13,000 & 0 & 0.00 & [0.000, 0.028] \\
control & none & 13,000 & 0 & 0.00 & [0.000, 0.028] \\
malicious & amount & 13,000 & 13,000 & 100.00 & [99.972, 100.000] \\
malicious & authorization & 13,000 & 13,000 & 100.00 & [99.972, 100.000] \\
malicious & beneficiary & 13,000 & 13,000 & 100.00 & [99.972, 100.000] \\
malicious & date & 13,000 & 13,000 & 100.00 & [99.972, 100.000] \\
malicious & document\_deleted & 13,000 & 13,000 & 100.00 & [99.972, 100.000] \\
malicious & invoice\_number & 13,000 & 13,000 & 100.00 & [99.972, 100.000] \\
malicious & ledger\_entry & 13,000 & 13,000 & 100.00 & [99.972, 100.000] \\
malicious & record\_substitution & 13,000 & 13,000 & 100.00 & [99.972, 100.000] \\
malicious & single\_bit\_metadata & 13,000 & 13,000 & 100.00 & [99.972, 100.000] \\
malicious & supporting\_document & 13,000 & 13,000 & 100.00 & [99.972, 100.000] \\
malicious & vendor & 13,000 & 13,000 & 100.00 & [99.972, 100.000] \\
malicious & ALL (TDR) & 143,000 & 143,000 & 100.00 & [99.997, 100.000] \\
benign & ALL (FPR) & 26,000 & 0 & 0.00 & [0.000, 0.014] \\
\bottomrule
\end{tabular}
\end{table}

\begin{table}[t]
\centering
\scriptsize
\caption{Tamper detection by dataset and infrastructure (E6).}
\label{tab:tamperds}
\resizebox{\linewidth}{!}{%
\begin{tabular}{llrrrrrr}
\toprule
Dataset & Infrastructure & Malicious & TDR (\%) & TDR 95\% lower & Benign & FPR (\%) & FPR 95\% upper \\
\midrule
IBM AML & PoA\textsuperscript{2} emulation & 33,000 & 100.00 & 99.989 & 6,000 & 0.00 & 0.061 \\
MoMTSim & PoA\textsuperscript{2} emulation & 33,000 & 100.00 & 99.989 & 6,000 & 0.00 & 0.061 \\
MoMTSim & PureChain (live) & 33,000 & 100.00 & 99.989 & 6,000 & 0.00 & 0.061 \\
SAML-D & PoA\textsuperscript{2} emulation & 33,000 & 100.00 & 99.989 & 6,000 & 0.00 & 0.061 \\
SAML-D & PureChain (live) & 11,000 & 100.00 & 99.966 & 2,000 & 0.00 & 0.184 \\
\bottomrule
\end{tabular}}
\end{table}


\subsection{Insider tampering (E8)}
Table~\ref{tab:insider} and Figure~\ref{fig:F11_insider} show the decisive difference between a trusted database and
a blockchain registry. When the insider only alters evidence (S1), every infrastructure detects it because the stored
commitment no longer matches. When the insider also rewrites the stored commitment (S2) or inserts a backdated
replacement record (S3), the database accepts the writes and the audit passes. On the blockchain configurations the
rewrite is rejected by the contract (records are write-once) and a fabricated record carries the block timestamp of its
actual inclusion, so both attacks are prevented or detected.
\begin{table}[t]
\centering
\scriptsize
\caption{Privileged-insider tampering (E8): share of attacks prevented (write rejected) or detected by the audit, with Clopper--Pearson 95\% intervals.}
\label{tab:insider}
\begin{tabular}{llrrrrc}
\toprule
Attack & Infrastructure & Cases & Prevented & Detected & Undetected & Protected (\%) \\
\midrule
S1 alter evidence & PoA\textsuperscript{2} emulation & 700 & 0 & 700 & 0 & 100.0 [99.5, 100.0] \\
S1 alter evidence & Besu QBFT & 700 & 0 & 700 & 0 & 100.0 [99.5, 100.0] \\
S1 alter evidence & PostgreSQL & 700 & 0 & 700 & 0 & 100.0 [99.5, 100.0] \\
S1 alter evidence & Ethereum PoS & 200 & 0 & 200 & 0 & 100.0 [98.2, 100.0] \\
S2 alter evidence + rewrite commitment & PoA\textsuperscript{2} emulation & 700 & 700 & 0 & 0 & 100.0 [99.5, 100.0] \\
S2 alter evidence + rewrite commitment & Besu QBFT & 700 & 700 & 0 & 0 & 100.0 [99.5, 100.0] \\
S2 alter evidence + rewrite commitment & PostgreSQL & 700 & 0 & 0 & 700 & 0.0 [0.0, 0.5] \\
S2 alter evidence + rewrite commitment & Ethereum PoS & 200 & 200 & 0 & 0 & 100.0 [98.2, 100.0] \\
S3 fabricate + backdate record & PoA\textsuperscript{2} emulation & 700 & 0 & 700 & 0 & 100.0 [99.5, 100.0] \\
S3 fabricate + backdate record & Besu QBFT & 700 & 0 & 700 & 0 & 100.0 [99.5, 100.0] \\
S3 fabricate + backdate record & PostgreSQL & 700 & 0 & 0 & 700 & 0.0 [0.0, 0.5] \\
S3 fabricate + backdate record & Ethereum PoS & 200 & 0 & 200 & 0 & 100.0 [98.2, 100.0] \\
\bottomrule
\end{tabular}
\end{table}

\figif{F11_insider}{0.62\linewidth}{Share of privileged-insider attacks prevented or detected (E8).}

\subsection{Batched Merkle commitments (E9)}
Committing $k$ records under one Merkle root divides the number of blockchain transactions and the gas per record by
$k$, at the cost of the time a record waits for its batch to close (at most 1\,s). Table~\ref{tab:batching} and
Figure~\ref{fig:F10_batching} report throughput and end-to-end latency as a function of $k$.
\begin{table}[t]
\centering
\scriptsize
\caption{Batched Merkle commitments (E9): records per commitment $k$ versus throughput, end-to-end inclusion latency and finality latency (mean over trials, 95\% CI).}
\label{tab:batching}
\resizebox{\linewidth}{!}{%
\begin{tabular}{lrrrcccr}
\toprule
Infrastructure & Offered & $k$ & $n$ & Throughput (rec/s) & Arrival$\rightarrow$incl. P50 (s) & Arrival$\rightarrow$final P50 (s) & Gas/record \\
\midrule
PoA\textsuperscript{2} emulation & 50 & 1 & 5 & 48.22 [47.44, 49.00] & 1.55 [1.53, 1.58] & 7.55 [7.53, 7.56] & 153,376 \\
PoA\textsuperscript{2} emulation & 50 & 10 & 5 & 48.15 [47.00, 49.30] & 1.62 [1.60, 1.64] & 7.62 [7.60, 7.64] & 15,932 \\
PoA\textsuperscript{2} emulation & 50 & 25 & 5 & 48.13 [47.30, 48.96] & 1.77 [1.76, 1.78] & 7.77 [7.76, 7.78] & 6,373 \\
PoA\textsuperscript{2} emulation & 50 & 50 & 5 & 47.92 [46.69, 49.15] & 2.01 [1.92, 2.10] & 8.01 [7.93, 8.10] & 3,377 \\
PoA\textsuperscript{2} emulation & 50 & 100 & 5 & 47.65 [46.64, 48.67] & 2.06 [1.97, 2.14] & 8.05 [7.97, 8.14] & 3,186 \\
PoA\textsuperscript{2} emulation & 200 & 1 & 5 & 183.24 [178.48, 187.99] & 2.10 [1.98, 2.22] & 8.10 [7.95, 8.24] & 153,377 \\
PoA\textsuperscript{2} emulation & 200 & 10 & 5 & 189.33 [185.18, 193.49] & 1.56 [1.55, 1.57] & 7.56 [7.55, 7.57] & 15,931 \\
PoA\textsuperscript{2} emulation & 200 & 25 & 5 & 191.41 [188.24, 194.58] & 1.60 [1.58, 1.61] & 7.60 [7.58, 7.61] & 6,372 \\
PoA\textsuperscript{2} emulation & 200 & 50 & 5 & 190.47 [186.45, 194.49] & 1.66 [1.63, 1.69] & 7.66 [7.63, 7.68] & 3,186 \\
PoA\textsuperscript{2} emulation & 200 & 100 & 5 & 190.08 [187.48, 192.67] & 1.77 [1.70, 1.83] & 7.77 [7.70, 7.84] & 1,593 \\
Besu QBFT & 50 & 1 & 5 & 43.83 [36.00, 51.67] & 7.25 [0.00, 17.43] & 7.25 [0.00, 17.43] & 153,378 \\
Besu QBFT & 50 & 10 & 5 & 48.94 [47.82, 50.07] & 0.69 [0.68, 0.71] & 0.69 [0.68, 0.71] & 15,932 \\
Besu QBFT & 50 & 25 & 5 & 48.97 [47.99, 49.95] & 0.82 [0.80, 0.85] & 0.82 [0.80, 0.85] & 6,372 \\
Besu QBFT & 50 & 50 & 5 & 48.56 [47.30, 49.82] & 1.02 [0.92, 1.12] & 1.02 [0.92, 1.12] & 3,377 \\
Besu QBFT & 50 & 100 & 5 & 48.33 [47.13, 49.53] & 1.11 [1.09, 1.13] & 1.11 [1.09, 1.13] & 3,186 \\
Besu QBFT & 200 & 1 & 5 & 51.79 [49.55, 54.04] & 50.71 [47.26, 54.16] & 50.71 [47.26, 54.16] & 153,377 \\
Besu QBFT & 200 & 10 & 5 & 194.97 [190.41, 199.54] & 0.87 [0.84, 0.89] & 0.87 [0.84, 0.89] & 15,931 \\
Besu QBFT & 200 & 25 & 5 & 195.59 [191.46, 199.73] & 0.70 [0.69, 0.71] & 0.70 [0.69, 0.71] & 6,373 \\
Besu QBFT & 200 & 50 & 5 & 195.49 [190.80, 200.18] & 0.72 [0.71, 0.74] & 0.72 [0.71, 0.74] & 3,186 \\
Besu QBFT & 200 & 100 & 5 & 196.71 [192.42, 200.99] & 0.83 [0.75, 0.91] & 0.83 [0.75, 0.91] & 1,593 \\
PostgreSQL & 50 & 1 & 5 & 49.41 [48.37, 50.44] & 0.00 [0.00, 0.00] & 0.00 [0.00, 0.00] & -- \\
PostgreSQL & 50 & 10 & 5 & 49.40 [48.37, 50.44] & 0.09 [0.09, 0.09] & 0.09 [0.09, 0.09] & -- \\
PostgreSQL & 50 & 25 & 5 & 49.40 [48.37, 50.44] & 0.24 [0.23, 0.25] & 0.24 [0.23, 0.25] & -- \\
PostgreSQL & 50 & 50 & 5 & 48.95 [47.79, 50.10] & 0.49 [0.47, 0.50] & 0.49 [0.47, 0.50] & -- \\
PostgreSQL & 50 & 100 & 5 & 48.97 [47.78, 50.15] & 0.52 [0.51, 0.53] & 0.52 [0.51, 0.53] & -- \\
PostgreSQL & 200 & 1 & 5 & 199.60 [196.12, 203.07] & 0.00 [0.00, 0.01] & 0.00 [0.00, 0.01] & -- \\
PostgreSQL & 200 & 10 & 5 & 199.59 [196.12, 203.06] & 0.03 [0.03, 0.03] & 0.03 [0.03, 0.03] & -- \\
PostgreSQL & 200 & 25 & 5 & 199.58 [196.11, 203.05] & 0.06 [0.06, 0.07] & 0.06 [0.06, 0.07] & -- \\
PostgreSQL & 200 & 50 & 5 & 199.56 [196.09, 203.03] & 0.13 [0.13, 0.13] & 0.13 [0.13, 0.13] & -- \\
PostgreSQL & 200 & 100 & 5 & 199.54 [196.07, 203.01] & 0.26 [0.26, 0.27] & 0.26 [0.26, 0.27] & -- \\
Ethereum PoS & 50 & 1 & 2 & 15.36 [15.35, 15.38] & 67.83 [62.79, 72.88] & 1,053 [439.68, 1,666] & 153,376 \\
Ethereum PoS & 50 & 50 & 2 & 44.92 [44.65, 45.18] & 8.72 [8.31, 9.12] & 1,105 [1,101, 1,108] & 3,375 \\
\bottomrule
\end{tabular}}
\end{table}

\figif{F10_batching}{\linewidth}{Batched commitments: throughput, arrival-to-finality latency and gas per record (E9).}

\subsection{Cross-dataset replication (E7)}
Table~\ref{tab:datasets} repeats the core workloads on the three datasets. The Kruskal--Wallis $p$-value tests whether
the dataset changes the median inclusion latency of an infrastructure; a large $p$ indicates that the results do not
depend on the workload source.
\begin{table}[t]
\centering
\scriptsize
\caption{Cross-dataset replication (E1 on MoMTSim, E7 on SAML-D and IBM AML): mean over trials with 95\% CI. $p$: Kruskal--Wallis test of a dataset effect on median inclusion latency.}
\label{tab:datasets}
\resizebox{\linewidth}{!}{%
\begin{tabular}{lllrcccrr}
\toprule
Workload & Infrastructure & Dataset & $n$ & Throughput (rec/s) & Incl. P50 (s) & Final P50 (s) & Verified (\%) & $p$ \\
\midrule
W1 & Besu QBFT & MoMTSim & 30 & 45.69 [44.25, 47.13] & 1.78 [0.97, 2.60] & 1.78 [0.97, 2.60] & 100.0 & 0.273 \\
 &  & SAML-D & 10 & 42.98 [38.81, 47.15] & 2.96 [0.50, 5.42] & 2.96 [0.50, 5.42] & 100.0 &  \\
 &  & IBM AML & 10 & 45.87 [45.01, 46.73] & 1.34 [1.28, 1.40] & 1.34 [1.28, 1.40] & 100.0 &  \\
W1 & PostgreSQL & MoMTSim & 30 & 49.58 [49.05, 50.11] & 0.00 [0.00, 0.00] & 0.00 [0.00, 0.00] & 100.0 & 0.000 \\
 &  & SAML-D & 10 & 49.28 [48.39, 50.16] & 0.00 [0.00, 0.00] & 0.00 [0.00, 0.00] & 100.0 &  \\
 &  & IBM AML & 10 & 49.28 [48.39, 50.16] & 0.00 [0.00, 0.00] & 0.00 [0.00, 0.00] & 100.0 &  \\
W1 & Ethereum PoS & MoMTSim & 5 & 14.71 [14.61, 14.80] & 22.54 [21.96, 23.12] & 1,126 [1,125, 1,127] & 100.0 & -- \\
W1 & PoA\textsuperscript{2} emulation & MoMTSim & 30 & 45.98 [45.48, 46.48] & 1.55 [1.54, 1.55] & 7.54 [7.53, 7.55] & 100.0 & 0.759 \\
 &  & SAML-D & 10 & 45.87 [45.17, 46.57] & 1.55 [1.54, 1.56] & 7.53 [7.50, 7.56] & 100.0 &  \\
 &  & IBM AML & 10 & 45.70 [44.97, 46.43] & 1.55 [1.53, 1.56] & 7.54 [7.53, 7.55] & 100.0 &  \\
W2 & Besu QBFT & MoMTSim & 10 & 49.69 [49.40, 49.98] & 1.28 [1.14, 1.41] & 1.28 [1.14, 1.41] & 100.0 & 0.004 \\
 &  & SAML-D & 5 & 49.68 [49.14, 50.21] & 1.33 [1.30, 1.37] & 1.33 [1.30, 1.37] & 100.0 &  \\
 &  & IBM AML & 5 & 49.56 [49.00, 50.13] & 1.47 [1.35, 1.59] & 1.47 [1.35, 1.59] & 100.0 &  \\
W2 & PostgreSQL & MoMTSim & 10 & 50.01 [49.71, 50.30] & 0.00 [0.00, 0.00] & 0.00 [0.00, 0.00] & 100.0 & 0.005 \\
 &  & SAML-D & 5 & 50.04 [49.48, 50.60] & 0.00 [0.00, 0.00] & 0.00 [0.00, 0.00] & 100.0 &  \\
 &  & IBM AML & 5 & 50.04 [49.48, 50.60] & 0.00 [0.00, 0.00] & 0.00 [0.00, 0.00] & 100.0 &  \\
W2 & Ethereum PoS & MoMTSim & 2 & 16.05 [15.97, 16.13] & 213.21 [210.02, 216.40] & 1,103 [1,098, 1,107] & 100.0 & -- \\
W2 & PoA\textsuperscript{2} emulation & MoMTSim & 10 & 49.63 [49.30, 49.97] & 1.56 [1.55, 1.56] & 7.55 [7.55, 7.56] & 100.0 & 0.742 \\
 &  & SAML-D & 5 & 49.74 [49.15, 50.32] & 1.55 [1.55, 1.56] & 7.55 [7.55, 7.56] & 100.0 &  \\
 &  & IBM AML & 5 & 49.64 [49.12, 50.16] & 1.55 [1.55, 1.56] & 7.55 [7.54, 7.56] & 100.0 &  \\
W5 & Besu QBFT & MoMTSim & 10 & 49.59 [49.28, 49.90] & 2.89 [0.00, 6.37] & 2.89 [0.00, 6.37] & 100.0 & 0.089 \\
 &  & SAML-D & 3 & 49.74 [49.45, 50.04] & 1.37 [1.30, 1.44] & 1.37 [1.30, 1.44] & 100.0 &  \\
 &  & IBM AML & 3 & 49.66 [49.01, 50.32] & 1.59 [1.42, 1.76] & 1.59 [1.42, 1.76] & 100.0 &  \\
W5 & PostgreSQL & MoMTSim & 10 & 50.01 [49.71, 50.30] & 0.00 [0.00, 0.00] & 0.00 [0.00, 0.00] & 100.0 & 0.349 \\
 &  & SAML-D & 3 & 50.15 [49.74, 50.56] & 0.00 [0.00, 0.01] & 0.00 [0.00, 0.01] & 100.0 &  \\
 &  & IBM AML & 3 & 50.15 [49.74, 50.56] & 0.00 [0.00, 0.01] & 0.00 [0.00, 0.01] & 100.0 &  \\
W5 & PoA\textsuperscript{2} emulation & MoMTSim & 10 & 49.64 [49.36, 49.92] & 1.55 [1.55, 1.56] & 7.55 [7.55, 7.56] & 100.0 & 0.147 \\
 &  & SAML-D & 3 & 49.79 [49.10, 50.49] & 1.56 [1.55, 1.56] & 7.55 [7.55, 7.56] & 100.0 &  \\
 &  & IBM AML & 3 & 49.71 [49.28, 50.14] & 1.56 [1.55, 1.57] & 7.56 [7.55, 7.57] & 100.0 &  \\
\bottomrule
\end{tabular}}
\end{table}


\subsection{Latency, finality and throughput (E1)}
At the common offered load of 50 records/s (W2, 10{,}000 records), the \poa{} emulation, Besu and PostgreSQL all
sustain the offered load (\TpsWtwoEmu, \TpsWtwoBesu{} and \TpsWtwoDB{} records/s). Median inclusion latency is
\InclWtwoEmu\,s for \poa{}, \InclWtwoBesu\,s for Besu QBFT and below 5\,ms for PostgreSQL. The Ethereum PoS
configuration saturates at \TpsWtwoPos{} records/s, so its inclusion latency grows to \InclWtwoPos\,s, which is
\InclRatioPos{} times the \poa{} value. Finality separates the protocols more sharply
(Figure~\ref{fig:F9_finality_ladder}): QBFT is final on inclusion (\FinalWtwoBesu\,s), \poa{} reaches its
\FinalityDepth-confirmation finality in \FinalWtwoEmu\,s, and PoS reaches a finalized checkpoint only after
\FinalWtwoPos\,s, a factor of \FinalRatioPos{} slower than \poa. Workloads W1--W5 show the same ordering. For \poa{}, throughput
and latency are essentially independent of workload volume up to 100{,}000 records, whereas the Besu latency varies
more between workloads and trials (Table~\ref{tab:e1}, Figure~\ref{fig:F3_workloads}).
\begin{table}[t]
\centering
\scriptsize
\caption{Controlled workloads (E1), offered load 50 records/s.}
\label{tab:e1}
\resizebox{\linewidth}{!}{%
\begin{tabular}{llrrcccccr}
\toprule
Workload & Infrastructure & $n$ & Success (\%) & Throughput (rec/s) & Incl. P50 (s) & Incl. P95 (s) & Final P50 (s) & J/record & \textmu USD/record \\
\midrule
W1 & Besu QBFT & 30 & 100.0 & 45.69 [44.25, 47.13] & 1.78 [0.97, 2.60] & 2.75 [1.28, 4.22] & 1.78 [0.97, 2.60] & 2.226 & 36.01 \\
W1 & PostgreSQL & 30 & 100.0 & 49.58 [49.05, 50.11] & 0.00 [0.00, 0.00] & 0.00 [0.00, 0.00] & 0.00 [0.00, 0.00] & 0.510 & 25.69 \\
W1 & Ethereum PoS & 5 & 100.0 & 14.71 [14.61, 14.80] & 22.54 [21.96, 23.12] & 39.65 [38.99, 40.31] & 1,126 [1,125, 1,127] & 6.814 & 2,300,681 \\
W1 & PoA\textsuperscript{2} emulation & 30 & 100.0 & 45.98 [45.48, 46.48] & 1.55 [1.54, 1.55] & 2.00 [2.00, 2.01] & 7.54 [7.53, 7.55] & 2.188 & 36.09 \\
W2 & Besu QBFT & 10 & 100.0 & 49.69 [49.40, 49.98] & 1.28 [1.14, 1.41] & 2.57 [0.71, 4.42] & 1.28 [1.14, 1.41] & 2.025 & 35.87 \\
W2 & PostgreSQL & 10 & 100.0 & 50.01 [49.71, 50.30] & 0.00 [0.00, 0.00] & 0.00 [0.00, 0.00] & 0.00 [0.00, 0.00] & 0.505 & 25.69 \\
W2 & Ethereum PoS & 2 & 100.0 & 16.05 [15.97, 16.13] & 213.21 [210.02, 216.40] & 398.73 [389.21, 408.25] & 1,103 [1,098, 1,107] & 6.249 & 2,300,688 \\
W2 & PoA\textsuperscript{2} emulation & 10 & 100.0 & 49.63 [49.30, 49.97] & 1.56 [1.55, 1.56] & 2.01 [2.00, 2.01] & 7.55 [7.55, 7.56] & 2.026 & 35.93 \\
W3 & Besu QBFT & 5 & 100.0 & 50.00 [49.67, 50.34] & 1.34 [1.29, 1.39] & 3.85 [0.00, 8.86] & 1.34 [1.29, 1.39] & 2.013 & 213.91 \\
W3 & PostgreSQL & 5 & 100.0 & 50.07 [49.73, 50.41] & 0.00 [0.00, 0.00] & 0.01 [0.01, 0.01] & 0.00 [0.00, 0.00] & 0.505 & 203.76 \\
W3 & PoA\textsuperscript{2} emulation & 5 & 100.0 & 49.98 [49.64, 50.33] & 1.56 [1.56, 1.56] & 2.01 [2.01, 2.02] & 7.56 [7.55, 7.56] & 2.012 & 213.98 \\
W4 & Besu QBFT & 3 & 100.0 & 50.10 [49.75, 50.45] & 1.39 [1.20, 1.58] & 2.72 [1.44, 4.00] & 1.39 [1.20, 1.58] & 2.009 & 213.90 \\
W4 & PostgreSQL & 3 & 100.0 & 50.14 [49.80, 50.48] & 0.01 [0.01, 0.01] & 0.01 [0.01, 0.01] & 0.01 [0.01, 0.01] & 0.505 & 203.76 \\
W4 & PoA\textsuperscript{2} emulation & 3 & 100.0 & 50.10 [49.75, 50.46] & 1.56 [1.55, 1.57] & 2.01 [2.00, 2.03] & 7.56 [7.55, 7.57] & 2.007 & 213.98 \\
W5 & Besu QBFT & 10 & 100.0 & 49.59 [49.28, 49.90] & 2.89 [0.00, 6.37] & 5.65 [0.00, 13.59] & 2.89 [0.00, 6.37] & 2.030 & 2,042 \\
W5 & PostgreSQL & 10 & 100.0 & 50.01 [49.71, 50.30] & 0.00 [0.00, 0.00] & 0.01 [0.01, 0.01] & 0.00 [0.00, 0.00] & 0.506 & 2,032 \\
W5 & PoA\textsuperscript{2} emulation & 10 & 100.0 & 49.64 [49.36, 49.92] & 1.55 [1.55, 1.56] & 2.01 [2.00, 2.01] & 7.55 [7.55, 7.56] & 2.026 & 2,042 \\
\bottomrule
\end{tabular}}
\end{table}

\figif{F1_latency_cdf}{0.55\linewidth}{Commit-latency distribution (E1, W2).}
\figif{F9_finality_ladder}{0.62\linewidth}{Inclusion versus confirmation-depth finality (E1, W2).}
\figif{F3_workloads}{\linewidth}{Latency and throughput across workloads W1--W5 (95\% CI).}

\subsection{Capacity under increasing load (E2)}
Increasing the offered load exposes the capacity of each configuration (Table~\ref{tab:e2},
Figure~\ref{fig:F2_offered_load_sweep}). The \poa{} emulation sustains \MaxTpsEmu{} records/s with per-record
commitments, \MaxTpsRatio{} times the \MaxTpsBesu{} records/s of Besu QBFT under the same per-validator CPU budget. At
200 records/s offered, \poa{} commits \TpsTwoHEmu{} records/s with a median inclusion latency of \LatTwoHEmu\,s,
whereas Besu is saturated at \TpsTwoHBesu{} records/s and its latency rises to \LatTwoHBesu\,s. PostgreSQL remains the
fastest system (\MaxTpsDB{} records/s unthrottled), as expected for a single trusted server without replication.
\begin{table}[t]
\centering
\scriptsize
\caption{Offered-load sweep (E2).}
\label{tab:e2}
\resizebox{\linewidth}{!}{%
\begin{tabular}{llrrcccccr}
\toprule
Offered (rec/s) & Infrastructure & $n$ & Success (\%) & Throughput (rec/s) & Incl. P50 (s) & Incl. P95 (s) & Final P50 (s) & J/record & \textmu USD/record \\
\midrule
10 & Besu QBFT & 5 & 100.0 & 9.78 [9.54, 10.03] & 0.70 [0.67, 0.73] & 1.17 [1.11, 1.22] & 0.70 [0.67, 0.73] & 10.259 & 37.16 \\
10 & PostgreSQL & 5 & 100.0 & 9.82 [9.57, 10.06] & 0.00 [0.00, 0.00] & 0.01 [0.01, 0.01] & 0.00 [0.00, 0.00] & 2.554 & 25.76 \\
10 & Ethereum PoS & 2 & 100.0 & 8.86 [7.36, 10.36] & 7.62 [5.26, 9.98] & 13.39 [12.25, 14.52] & 1,106 [1,093, 1,118] & 11.308 & 2,300,686 \\
10 & PoA\textsuperscript{2} emulation & 5 & 100.0 & 9.74 [9.51, 9.97] & 1.54 [1.53, 1.55] & 1.98 [1.97, 1.99] & 7.54 [7.53, 7.55] & 10.299 & 37.01 \\
25 & Besu QBFT & 5 & 100.0 & 24.27 [23.67, 24.87] & 0.87 [0.84, 0.91] & 1.35 [1.31, 1.39] & 0.87 [0.84, 0.91] & 4.141 & 36.27 \\
25 & PostgreSQL & 5 & 100.0 & 24.54 [23.94, 25.15] & 0.00 [0.00, 0.00] & 0.01 [0.00, 0.01] & 0.00 [0.00, 0.00] & 1.026 & 25.71 \\
25 & PoA\textsuperscript{2} emulation & 5 & 100.0 & 24.09 [23.43, 24.75] & 1.53 [1.52, 1.54] & 1.99 [1.98, 1.99] & 7.53 [7.52, 7.54] & 4.169 & 36.26 \\
50 & Besu QBFT & 5 & 100.0 & 47.46 [46.29, 48.64] & 1.25 [1.21, 1.29] & 1.80 [1.72, 1.88] & 1.25 [1.21, 1.29] & 2.121 & 35.92 \\
50 & PostgreSQL & 5 & 100.0 & 49.08 [47.87, 50.30] & 0.00 [0.00, 0.00] & 0.00 [0.00, 0.00] & 0.00 [0.00, 0.00] & 0.515 & 25.69 \\
50 & Ethereum PoS & 2 & 100.0 & 15.43 [14.41, 16.44] & 68.09 [67.19, 69.00] & 124.40 [120.91, 127.89] & 1,100 [1,080, 1,120] & 6.497 & 2,300,684 \\
50 & PoA\textsuperscript{2} emulation & 5 & 100.0 & 47.51 [46.16, 48.87] & 1.56 [1.54, 1.58] & 2.00 [1.99, 2.00] & 7.55 [7.53, 7.57] & 2.118 & 36.01 \\
100 & Besu QBFT & 5 & 100.0 & 58.12 [55.10, 61.14] & 11.34 [8.84, 13.84] & 25.78 [22.26, 29.30] & 11.34 [8.84, 13.84] & 1.736 & 36.12 \\
100 & PostgreSQL & 5 & 100.0 & 98.81 [96.74, 100.87] & 0.00 [0.00, 0.00] & 0.00 [0.00, 0.00] & 0.00 [0.00, 0.00] & 0.259 & 25.68 \\
100 & PoA\textsuperscript{2} emulation & 5 & 100.0 & 94.20 [92.44, 95.95] & 1.58 [1.57, 1.59] & 2.04 [2.03, 2.05] & 7.56 [7.55, 7.57] & 1.072 & 35.84 \\
200 & Besu QBFT & 5 & 100.0 & 57.90 [54.72, 61.08] & 46.21 [41.13, 51.29] & 73.82 [68.18, 79.47] & 46.21 [41.13, 51.29] & 1.745 & 36.52 \\
200 & PostgreSQL & 5 & 100.0 & 199.61 [196.13, 203.08] & 0.00 [0.00, 0.00] & 0.00 [0.00, 0.00] & 0.00 [0.00, 0.00] & 0.131 & 25.68 \\
200 & PoA\textsuperscript{2} emulation & 5 & 100.0 & 184.38 [183.91, 184.84] & 1.99 [1.79, 2.19] & 2.72 [2.34, 3.10] & 7.97 [7.76, 8.19] & 0.551 & 35.72 \\
400 & Besu QBFT & 5 & 100.0 & 54.60 [51.64, 57.55] & 110.64 [100.27, 121.01] & 185.20 [173.79, 196.62] & 110.64 [100.27, 121.01] & 1.852 & 36.87 \\
400 & PostgreSQL & 5 & 100.0 & 401.41 [398.17, 404.66] & 0.00 [0.00, 0.00] & 0.01 [0.00, 0.01] & 0.00 [0.00, 0.00] & 0.068 & 25.67 \\
400 & PoA\textsuperscript{2} emulation & 5 & 100.0 & 191.14 [190.50, 191.78] & 17.33 [16.94, 17.73] & 31.38 [31.05, 31.71] & 23.13 [22.76, 23.50] & 0.532 & 35.79 \\
max & Besu QBFT & 5 & 100.0 & 56.69 [54.85, 58.52] & 51.74 [49.13, 54.34] & 85.70 [82.25, 89.16] & 51.74 [49.13, 54.34] & 1.781 & 36.54 \\
max & PostgreSQL & 5 & 100.0 & 1,378 [1,344, 1,411] & 1.81 [1.76, 1.86] & 3.42 [3.34, 3.50] & 1.81 [1.76, 1.86] & 0.025 & 25.65 \\
max & PoA\textsuperscript{2} emulation & 5 & 100.0 & 186.77 [184.21, 189.34] & 13.83 [13.42, 14.24] & 25.88 [25.53, 26.22] & 19.81 [19.44, 20.18] & 0.545 & 35.84 \\
\bottomrule
\end{tabular}}
\end{table}

\figif{F2_offered_load_sweep}{\linewidth}{Throughput and latency versus offered load (E2).}

\subsection{Validator scaling (E5)}
With every validator limited to one CPU, the \poa{} emulation keeps its throughput at 200 records/s offered as the
validator set grows (\VTpsEmuFour, \VTpsEmuSeven{} and \VTpsEmuTen{} records/s for 4, 7 and 10 signers). Besu QBFT
commits \VTpsBesuFour{} and \VTpsBesuSeven{} records/s with 4 and 7 validators; with 10 validators its BFT event queues
overflow and the network loses liveness, committing only \BesuTenSuccess\% of the records offered at 50 records/s
(Table~\ref{tab:e5}, Figure~\ref{fig:F6_validators}). A fresh confirmation run with 10 validators completed but ran at
less than half the offered rate. The all-to-all message pattern of BFT consensus grows quadratically with the
validator count, while Clique-style sealing requires one signature per block.
\begin{table}[t]
\centering
\scriptsize
\caption{Validator-count sweep (E5); every validator is limited to 1 CPU. Besu QBFT with 10 validators lost liveness in 3 of 4 trials (BFT event-queue overflow).}
\label{tab:e5}
\resizebox{\linewidth}{!}{%
\begin{tabular}{llrrcccccr}
\toprule
Validators @ offered & Infrastructure & $n$ & Success (\%) & Throughput (rec/s) & Incl. P50 (s) & Incl. P95 (s) & Final P50 (s) & J/record & \textmu USD/record \\
\midrule
v10 @ 50 & Besu QBFT & 3 & 4.0 & 3.34 [0.00, 17.74] & 14.12 & 31.79 & 14.12 & 24.987 & 536.70 \\
v10 @ 50 & PoA\textsuperscript{2} emulation & 5 & 100.0 & 48.20 [47.25, 49.15] & 1.58 [1.57, 1.59] & 2.21 [1.95, 2.47] & 7.57 [7.55, 7.58] & 5.201 & 52.10 \\
v10 @ 200 & PoA\textsuperscript{2} emulation & 5 & 100.0 & 174.64 [170.51, 178.76] & 2.03 [1.79, 2.27] & 2.66 [2.34, 2.98] & 7.89 [7.70, 8.08] & 1.441 & 51.77 \\
v10 (confirm) @ 50 & Besu QBFT & 1 & 100.0 & 22.53 & 32.61 & 46.99 & 32.61 & 11.119 & 53.01 \\
v4 @ 50 & Besu QBFT & 5 & 100.0 & 46.33 [41.08, 51.57] & 3.92 [0.00, 11.15] & 7.12 [0.00, 21.58] & 3.92 [0.00, 11.15] & 2.189 & 35.93 \\
v4 @ 50 & PoA\textsuperscript{2} emulation & 5 & 100.0 & 48.22 [47.22, 49.23] & 1.56 [1.55, 1.57] & 2.00 [2.00, 2.01] & 7.56 [7.55, 7.57] & 2.085 & 35.97 \\
v4 @ 200 & Besu QBFT & 5 & 100.0 & 60.34 [57.82, 62.85] & 22.62 [20.72, 24.51] & 36.69 [34.22, 39.17] & 22.62 [20.72, 24.51] & 1.673 & 36.31 \\
v4 @ 200 & PoA\textsuperscript{2} emulation & 5 & 100.0 & 176.63 [175.08, 178.19] & 1.79 [1.68, 1.90] & 2.30 [2.17, 2.42] & 7.74 [7.64, 7.84] & 0.574 & 35.76 \\
v7 @ 50 & Besu QBFT & 5 & 100.0 & 44.88 [38.06, 51.70] & 6.70 [0.00, 18.44] & 11.68 [0.00, 32.35] & 6.70 [0.00, 18.44] & 3.970 & 44.04 \\
v7 @ 50 & PoA\textsuperscript{2} emulation & 5 & 100.0 & 48.10 [47.18, 49.02] & 1.56 [1.52, 1.59] & 2.04 [2.03, 2.06] & 7.55 [7.52, 7.59] & 3.651 & 44.03 \\
v7 @ 200 & Besu QBFT & 5 & 100.0 & 54.31 [51.92, 56.70] & 24.82 [23.89, 25.75] & 41.60 [39.40, 43.81] & 24.82 [23.89, 25.75] & 3.241 & 44.46 \\
v7 @ 200 & PoA\textsuperscript{2} emulation & 5 & 100.0 & 175.46 [171.04, 179.88] & 1.93 [1.75, 2.10] & 2.52 [2.33, 2.72] & 7.87 [7.71, 8.04] & 1.007 & 43.77 \\
\bottomrule
\end{tabular}}
\end{table}

\figif{F6_validators}{\linewidth}{Throughput and latency versus validator count (E5).}

\subsection{Evidence size and complexity (E3, E4)}
Because only the 32-byte commitment is recorded, commit latency does not depend on the size or complexity of the
evidence (Tables~\ref{tab:e3} and~\ref{tab:e4}). Growth in evidence size and complexity is confined to client-side
processing (generation, hashing and off-chain storage), which rises from about 1\,ms per record without documents to
tens of milliseconds for 5\,MB of supporting documents (Figure~\ref{fig:F4_evidence_size}).
\begin{table}[t]
\centering
\scriptsize
\caption{Evidence-size sweep (E3).}
\label{tab:e3}
\resizebox{\linewidth}{!}{%
\begin{tabular}{llrrcccccr}
\toprule
Evidence (bytes) & Infrastructure & $n$ & Success (\%) & Throughput (rec/s) & Incl. P50 (s) & Incl. P95 (s) & Final P50 (s) & J/record & \textmu USD/record \\
\midrule
none & Besu QBFT & 5 & 100.0 & 40.37 [32.11, 48.64] & 4.18 [0.00, 9.36] & 7.27 [0.00, 16.49] & 4.18 [0.00, 9.36] & 2.555 & 15.23 \\
none & PostgreSQL & 5 & 100.0 & 48.93 [47.71, 50.16] & 0.00 [0.00, 0.00] & 0.01 [0.00, 0.01] & 0.00 [0.00, 0.00] & 0.518 & 4.87 \\
none & PoA\textsuperscript{2} emulation & 5 & 100.0 & 45.26 [44.24, 46.27] & 1.54 [1.53, 1.55] & 2.00 [1.99, 2.00] & 7.53 [7.52, 7.55] & 2.223 & 15.25 \\
10\,KB & Besu QBFT & 5 & 100.0 & 45.73 [44.13, 47.34] & 1.37 [1.27, 1.47] & 2.09 [1.84, 2.34] & 1.37 [1.27, 1.47] & 2.202 & 36.01 \\
10\,KB & PostgreSQL & 5 & 100.0 & 48.93 [47.71, 50.16] & 0.00 [0.00, 0.00] & 0.01 [0.00, 0.01] & 0.00 [0.00, 0.00] & 0.516 & 25.70 \\
10\,KB & PoA\textsuperscript{2} emulation & 5 & 100.0 & 45.67 [44.30, 47.04] & 1.54 [1.50, 1.57] & 2.00 [1.99, 2.01] & 7.53 [7.50, 7.56] & 2.203 & 36.11 \\
100\,KB & Besu QBFT & 5 & 100.0 & 45.94 [44.68, 47.20] & 1.29 [1.21, 1.37] & 1.87 [1.82, 1.92] & 1.29 [1.21, 1.37] & 2.192 & 214.07 \\
100\,KB & PostgreSQL & 5 & 100.0 & 48.93 [47.71, 50.16] & 0.00 [0.00, 0.00] & 0.00 [0.00, 0.01] & 0.00 [0.00, 0.00] & 0.517 & 203.75 \\
100\,KB & PoA\textsuperscript{2} emulation & 5 & 100.0 & 45.49 [44.21, 46.76] & 1.54 [1.52, 1.56] & 2.00 [1.99, 2.01] & 7.54 [7.52, 7.55] & 2.211 & 214.18 \\
1\,MB & Besu QBFT & 5 & 100.0 & 45.78 [44.09, 47.48] & 1.29 [1.22, 1.37] & 1.91 [1.72, 2.10] & 1.29 [1.22, 1.37] & 2.200 & 2,042 \\
1\,MB & PostgreSQL & 5 & 100.0 & 48.93 [47.71, 50.16] & 0.00 [0.00, 0.00] & 0.01 [0.01, 0.01] & 0.00 [0.00, 0.00] & 0.517 & 2,032 \\
1\,MB & PoA\textsuperscript{2} emulation & 5 & 100.0 & 45.09 [44.10, 46.07] & 1.55 [1.53, 1.57] & 2.01 [1.99, 2.02] & 7.55 [7.53, 7.56] & 2.231 & 2,042 \\
5\,MB & Besu QBFT & 5 & 100.0 & 45.90 [44.65, 47.15] & 1.31 [1.27, 1.35] & 1.89 [1.76, 2.02] & 1.31 [1.27, 1.35] & 2.194 & 10,149 \\
5\,MB & PostgreSQL & 5 & 100.0 & 48.93 [47.71, 50.15] & 0.01 [0.01, 0.01] & 0.01 [0.01, 0.01] & 0.01 [0.01, 0.01] & 0.517 & 10,135 \\
5\,MB & PoA\textsuperscript{2} emulation & 5 & 100.0 & 45.63 [44.71, 46.54] & 1.55 [1.53, 1.57] & 2.00 [1.99, 2.02] & 7.54 [7.52, 7.55] & 2.205 & 10,146 \\
\bottomrule
\end{tabular}}
\end{table}

\figif{F4_evidence_size}{\linewidth}{Evidence processing and commit latency versus evidence size (E3).}
\begin{table}[t]
\centering
\scriptsize
\caption{Evidence-complexity sweep (E4).}
\label{tab:e4}
\resizebox{\linewidth}{!}{%
\begin{tabular}{llrrcccccr}
\toprule
Level & Infrastructure & $n$ & Success (\%) & Throughput (rec/s) & Incl. P50 (s) & Incl. P95 (s) & Final P50 (s) & J/record & \textmu USD/record \\
\midrule
1 & Besu QBFT & 5 & 100.0 & 39.61 [28.77, 50.45] & 4.69 [0.00, 10.68] & 8.11 [0.00, 19.45] & 4.69 [0.00, 10.68] & 2.682 & 13.16 \\
1 & PostgreSQL & 5 & 100.0 & 48.93 [47.71, 50.16] & 0.00 [0.00, 0.00] & 0.01 [0.00, 0.01] & 0.00 [0.00, 0.00] & 0.518 & 2.77 \\
1 & PoA\textsuperscript{2} emulation & 5 & 100.0 & 45.13 [43.94, 46.33] & 1.56 [1.54, 1.59] & 2.01 [2.00, 2.01] & 7.56 [7.53, 7.58] & 2.228 & 13.15 \\
2 & Besu QBFT & 5 & 100.0 & 45.82 [44.48, 47.16] & 1.36 [1.13, 1.59] & 2.17 [1.50, 2.84] & 1.36 [1.13, 1.59] & 2.198 & 212.29 \\
2 & PostgreSQL & 5 & 100.0 & 48.93 [47.71, 50.16] & 0.00 [0.00, 0.00] & 0.01 [0.01, 0.01] & 0.00 [0.00, 0.00] & 0.517 & 201.97 \\
2 & PoA\textsuperscript{2} emulation & 5 & 100.0 & 45.07 [43.93, 46.20] & 1.55 [1.53, 1.58] & 2.00 [1.99, 2.01] & 7.54 [7.52, 7.57] & 2.232 & 212.39 \\
3 & Besu QBFT & 5 & 100.0 & 46.21 [45.06, 47.36] & 1.26 [1.17, 1.34] & 1.81 [1.75, 1.87] & 1.26 [1.17, 1.34] & 2.179 & 213.17 \\
3 & PostgreSQL & 5 & 100.0 & 48.93 [47.71, 50.15] & 0.00 [0.00, 0.00] & 0.01 [0.00, 0.01] & 0.00 [0.00, 0.00] & 0.517 & 202.86 \\
3 & PoA\textsuperscript{2} emulation & 5 & 100.0 & 45.58 [44.50, 46.66] & 1.56 [1.53, 1.59] & 2.00 [1.99, 2.02] & 7.55 [7.53, 7.57] & 2.208 & 213.27 \\
4 & Besu QBFT & 5 & 100.0 & 45.94 [44.23, 47.64] & 1.30 [1.26, 1.35] & 1.90 [1.81, 1.99] & 1.30 [1.26, 1.35] & 2.193 & 214.06 \\
4 & PostgreSQL & 5 & 100.0 & 48.93 [47.71, 50.16] & 0.00 [0.00, 0.00] & 0.00 [0.00, 0.01] & 0.00 [0.00, 0.00] & 0.517 & 203.75 \\
4 & PoA\textsuperscript{2} emulation & 5 & 100.0 & 45.58 [44.16, 46.99] & 1.54 [1.51, 1.57] & 2.00 [1.99, 2.02] & 7.53 [7.50, 7.56] & 2.207 & 214.16 \\
\bottomrule
\end{tabular}}
\end{table}

\figif{F5_complexity}{\linewidth}{Effect of evidence complexity (E4). The client code is identical for every
backend; the higher client cost during Besu trials reflects CPU contention from the co-located Java validators.}

\subsection{Efficiency: energy, resources and cost}
Under the attributed model every four-node blockchain uses \EattEmu\,J per record at 50 records/s, and the PoS
configuration \EattPos\,J (\EattRatioPos{} times more), because its lower throughput keeps the nodes busy longer for the
same number of records. The measured dynamic energy of the \poa{} emulation is \EdynEmu\,J per record, \EdynRatioBesu{}
times lower than Besu QBFT (\EdynBesu\,J) and below Ethereum PoS (\EdynPos\,J); part of this difference comes from the
client implementations (Go versus Java). Operational cost without fees is \CostEmu{} µUSD per record for \poa{},
identical to Besu and dominated by the 7-year storage of off-chain evidence. The public-chain configuration adds a gas
fee of \FeePos\,USD per record at the reference price, which PureChain's zero-gas design removes entirely
(Table~\ref{tab:efficiency}, Figure~\ref{fig:F7_energy_cost}). Continuous verification of a record takes
\VerifyEmu\,ms on \poa{}, against \VerifyBesu\,ms on Besu and \VerifyDB\,ms on PostgreSQL
(Figure~\ref{fig:F8_verification}).
\begin{table}[t]
\centering
\scriptsize
\caption{Efficiency metrics (MoMTSim workloads): CPU in cores used on average during the processing window, peak infrastructure memory, network I/O, storage overhead per record, energy and operational cost. n/m: not measurable (remote validators).}
\label{tab:efficiency}
\resizebox{\linewidth}{!}{%
\begin{tabular}{llrrrrrrrrr}
\toprule
Workload & Infrastructure & CPU client & CPU infra & Mem (MB) & Net (KB/rec) & On-chain (B/rec) & Off-chain (KB/rec) & E attr. (J) & E dyn. (J) & \textmu USD/rec \\
\midrule
W1 & Besu QBFT & 0.12 & 2.48 & 3,140 & 9.97 & 397 & 12.3 & 2.226 & 0.3125 & 36.01 \\
W1 & PostgreSQL & 0.05 & 0.02 & 80 & 0.90 & 182 & 12.3 & 0.510 & 0.0076 & 25.69 \\
W1 & Ethereum PoS & 0.04 & 0.30 & 835 & 47.99 & 402 & 12.3 & 6.814 & 0.1261 & 2,300,681 \\
W1 & PoA\textsuperscript{2} emulation & 0.10 & 0.30 & 1,207 & 10.43 & 398 & 12.3 & 2.188 & 0.0466 & 36.09 \\
W2 & Besu QBFT & 0.12 & 2.65 & 5,209 & 8.42 & 397 & 12.3 & 2.025 & 0.3016 & 35.87 \\
W2 & PostgreSQL & 0.05 & 0.02 & 87 & 0.94 & 182 & 12.3 & 0.505 & 0.0077 & 25.69 \\
W2 & Ethereum PoS & 0.06 & 0.45 & 1,012 & 56.65 & 402 & 12.3 & 6.249 & 0.1708 & 2,300,688 \\
W2 & PoA\textsuperscript{2} emulation & 0.10 & 0.36 & 2,726 & 8.69 & 398 & 12.3 & 2.026 & 0.0500 & 35.93 \\
W5 & Besu QBFT & 0.12 & 2.76 & 5,217 & 8.45 & 397 & 1,026 & 2.030 & 0.3151 & 2,042 \\
W5 & PostgreSQL & 0.05 & 0.02 & 105 & 0.94 & 182 & 1,026 & 0.506 & 0.0081 & 2,032 \\
W5 & PoA\textsuperscript{2} emulation & 0.10 & 0.36 & 2,597 & 8.68 & 398 & 1,026 & 2.026 & 0.0508 & 2,042 \\
\bottomrule
\end{tabular}}
\end{table}

\figif{F7_energy_cost}{\linewidth}{Energy and total cost per record (E1, W2), including gas fees.}
\figif{F8_verification}{0.62\linewidth}{Continuous verification time per record.}

% =====================================================================================
\clearpage
\section{Summary of findings}\label{sec:summary}
The research question asked whether a lightweight, zero-gas, \poa-based permissioned blockchain can provide
cryptographically verifiable audit evidence while reducing latency, energy consumption and operational cost.
\begin{itemize}
\item \textbf{Verifiable evidence.} Every committed record verified; \TDR\% of \TamperCases{} malicious modifications
were detected with no false positives, on all three datasets. Against a privileged insider, the blockchain registry
prevented or detected 100\% of commitment rewrites and backdated records, while the centralised database detected
none of them.
\item \textbf{Versus the conventional public-chain configuration.} \poa{} lowers inclusion latency by a factor of
\InclRatioPos, finality latency by a factor of \FinalRatioPos{} (\FinalWtwoEmu\,s versus \FinalWtwoPos\,s),
attributed energy by a factor of \EattRatioPos, and removes the gas fee of \FeePos\,USD per record.
\item \textbf{Versus the conventional permissioned blockchain.} \poa{} sustains \MaxTpsRatio{} times the per-record
capacity of Besu QBFT, keeps its throughput from 4 to 10 validators where QBFT loses liveness, and uses
\EdynRatioBesu{} times less measured dynamic energy. QBFT provides immediate finality, whereas \poa{} needs
\FinalityDepth{} confirmations (\FinalWtwoEmu\,s).
\item \textbf{Versus the centralised database.} The database is faster and about 30\% cheaper per record, but it
provides no protection against an insider who controls it; the blockchain overhead buys independent verifiability.
\item \textbf{Robustness.} The results do not depend on the workload dataset (smallest Kruskal--Wallis $p$ for
\poa{}: \KWminEmu) or on evidence size and complexity. Batched Merkle commitments cut gas per record from
\GasKone{} to \GasKhundred{} and let even the PoS configuration raise throughput from \PosTpsKone{} to
\PosTpsKfifty{} records/s, but they do not change its finality latency.
\end{itemize}

% =====================================================================================
\section{Threats to validity}\label{sec:validity}
\begin{itemize}
\item \textbf{Availability of the live network.} The public PureChain RPC endpoint stopped producing blocks during the
campaign (last block 1{,}807{,}036) while all four validators were still sealing immediately before the halt. The
controlled experiments were therefore completed on a calibrated \poa{} emulation; live measurements are reported
separately and will be extended when the endpoint recovers.
\item \textbf{Emulation fidelity.} The emulation reproduces validator count, block period, gas limit, fee-free
operation and Clique-style sealing, but not PureChain's modified fork rules, the contextual-association component of
\poa{}, the fork rate observed on the live network under load, or the wide-area network path.
\item \textbf{Live network stability.} Under sustained load of about 50 transactions/s the live network produced
competing blocks (\LiveReorgPct\% of observed blocks) and then halted; this is a property of the deployment at the
time of measurement and is reported, not excluded.
\item \textbf{Load-balanced endpoint.} The public endpoint alternates requests between two RPC nodes, so consecutive
nonces of one account can reach different nodes and be delayed. Live latencies are therefore conservative.
\item \textbf{Single host.} Local networks share one workstation; per-validator CPU quotas and backend isolation limit,
but do not remove, co-location effects.
\item \textbf{Energy model.} Remote validator energy cannot be metered; the attributed model uses a fixed node power and
is reported with the measured dynamic energy of the client and local containers.
\item \textbf{Synthetic evidence.} Audit evidence is generated synthetically because real audit evidence is
confidential; document contents do not affect hashing cost, which depends only on size.
\end{itemize}

\section*{Data and code availability}
Workload preprocessing, evidence generation, contract, load generator, analysis and this report are produced by the
FinVerity simulation framework; every table and figure is regenerated from the raw trial data with
\texttt{python -m analysis.report}.

\nocite{*}
\bibliographystyle{plain}
\bibliography{references}
\end{document}
